\documentclass[11pt]{article}

\usepackage[margin=1in]{geometry}
\usepackage{amsmath,amssymb,mathtools,bm}
\usepackage{booktabs}
\usepackage{microtype}
\usepackage[hidelinks]{hyperref}

\newcommand{\R}{\mathbb{R}}
\newcommand{\concat}{\mathbin{\Vert}}
\newcommand{\ind}{\mathbb{I}}
\newcommand{\main}{\mathrm{main}}
\newcommand{\rel}{\mathrm{rel}}
\DeclareMathOperator{\ReLU}{ReLU}
\DeclareMathOperator{\RMSNorm}{RMSNorm}
\DeclareMathOperator{\LayerNorm}{LayerNorm}
\DeclareMathOperator{\sigmoid}{sigmoid}
\DeclareMathOperator{\src}{src}
\DeclareMathOperator{\dst}{dst}
\DeclareMathOperator{\table}{table}

\title{Methodology: Relation-Conditioned Row-Graph Learning}
\author{}
\date{}

\begin{document}
\maketitle

\section{Methodology}
\label{sec:methodology}

\subsection{Problem Formulation and Architecture Overview}
\label{sec:problem-overview}

We consider binary prediction over a relational database.  The database is
represented at row level by
\begin{equation}
    \mathcal G^{\mathrm{DB}}
    =
    \bigl(\mathcal V^{\mathrm{DB}},\mathcal E^{\mathrm{FK}}\bigr),
    \label{eq:database-graph}
\end{equation}
where \(\mathcal V^{\mathrm{DB}}\) contains all database rows and
\(\mathcal E^{\mathrm{FK}}\) is an edge-index set containing the directed
child-to-parent row links induced by foreign-key/primary-key references.
Each edge \(e\in\mathcal E^{\mathrm{FK}}\) has endpoints \(\src(e)\) and
\(\dst(e)\); treating edges as indexed objects allows parallel FK links
between the same pair of rows.  Let \(\mathfrak T\) be the set of binary
prediction tasks, and let \(\kappa(g)\in\mathfrak T\) identify the task to
which prediction instance \(g\) belongs.  The instance is specified by a seed row
\(r_g^\star\in\mathcal V^{\mathrm{DB}}\), a target column \(c_g^\star\), and
a binary target \(y_g\in\{0,1\}\).  Let \(n_g^\star\) denote the target cell
located at \((r_g^\star,c_g^\star)\).  Its \emph{value} is masked and
excluded from cell encoding.
Its structural metadata---including the target column and semantic type,
seed table, and seed timestamp---remain available for graph construction,
candidate selection, and relative-time encoding.

Our model first samples a bounded relational context around \(r_g^\star\),
then constructs two query-specific graphs.  The \emph{main graph} contains
the sampled rows and two output-candidate nodes, whereas the auxiliary
\emph{relation graph} contains the foreign-key relation types observed in
the sampled context.  Two rounds of relation-graph propagation produce a
condition for each foreign-key link.  These conditions modulate six rounds
of typed message passing on the main graph.  Finally, a shared readout scores
the compatibility between the seed and each output candidate.  The base
variant omits relation-graph computation and sets every FK condition to
zero.

We use \(d\) for the hidden dimension, \(d_{\mathrm{text}}\) for the
dimension of precomputed text features, \(F\) for the number of temporal
frequencies, \(P\) for the number of relation-graph layers, and \(L\) for the
number of main-graph layers.  We retain \(d\) symbolically throughout the
derivation.  The classification configuration documented here uses
\begin{equation}
    d=256,
    \qquad
    d_{\mathrm{text}}=384,
    \qquad
    F=16,
    \qquad
    P=2,
    \qquad
    L=6.
    \label{eq:model-dimensions}
\end{equation}

\paragraph{Notation.}
We use \(g\) for a prediction instance, \(r\) for a row, \(n\) for a cell,
\(k\) for a candidate, and \(i,j\) for generic main-graph nodes.  Main-graph
edges and layers are indexed by \(e\) and \(\ell\), whereas relation-graph
nodes, edges, and layers are indexed by \(\rho\), \(\epsilon\), and \(p\),
respectively.  The maps \(\tau_g\) and \(\eta_g\) assign main-graph and
relation-graph edge types.  Bold lowercase symbols denote vectors, bold
uppercase symbols denote matrices, calligraphic symbols denote sets or
graphs, and unbold symbols denote scalars or indices.

All learnable parameters introduced below are shared across
\(\kappa\in\mathfrak T\) and query instances.  Consequently, no learnable matrix or bias has a
task or graph index.  Parameters remain distinct across semantic types,
edge types, and neural-network layers unless explicitly stated otherwise.


\subsection{Seed-Centered Relational Context Sampling}
\label{sec:context-sampling}

Applying the model to the complete database graph is generally
intractable.  For each instance \(g\), we therefore perform a bounded,
seed-centered, priority-based hybrid traversal of
\(\mathcal G^{\mathrm{DB}}\).  Starting from \(r_g^\star\), the sampler
traverses both foreign-to-primary and primary-to-foreign links, maintains a
visited set to avoid duplicate rows, and serializes the cells of every
newly visited row.  Foreign-to-primary expansions use a prioritized LIFO
frontier and may therefore follow a parent chain depth-first.  Once that
frontier is empty, the sampler selects reverse-FK neighbors from the
shallowest nonempty depth, randomly ordering neighbors within that depth.
The resulting cell sequence is denoted by
\begin{equation}
    \mathcal S_g
    =
    \operatorname{Sample}
    \bigl(
        \mathcal G^{\mathrm{DB}},r_g^\star;
        S_{\max},B_{\max}
    \bigr),
    \qquad
    |\mathcal S_g|\leq S_{\max},
    \label{eq:context-sample}
\end{equation}
where \(S_{\max}=2048\) is a \emph{cell} budget rather than a row budget.
Cells are appended row by row until this budget is exhausted; hence, the
last sampled row may be only partially represented.  Sequences shorter than
\(S_{\max}\) are padded.

To prevent a high-degree primary row from exhausting the context budget,
the sampler retains at most \(B_{\max}=256\) randomly selected database
children in a single primary-to-foreign expansion.  For forecasting
instances, a primary-to-foreign neighbor with an available timestamp is
discarded when it occurs after the seed timestamp.  Thus, this temporal
filter is applied specifically during reverse-FK expansion rather than as a
global filter over all rows.  Task-defined leakage columns are additionally
omitted from the seed and from rows whose stored optional timestamp equals
the seed's; under the implementation's equality check, this also includes
the case in which both timestamps are missing.

Let the cells supplied to the encoder be
\begin{equation}
    \mathcal C_g^{\mathrm{in}}
    =
    \left\{
        n\in\mathcal S_g
        \;\middle|\;
        n\text{ is not masked}
    \right\}.
    \label{eq:input-cell-set}
\end{equation}
The row nodes are precisely the rows containing at least one such cell:
\begin{equation}
    \mathcal R_g
    =
    \bigl\{
        r(n)\;\big|\;n\in\mathcal C_g^{\mathrm{in}}
    \bigr\},
    \label{eq:sampled-row-set}
\end{equation}
where \(r(n)\) and \(c(n)\) denote the row and column containing cell \(n\),
respectively.  Here
\(\mathcal S_g\) denotes the non-padding sampled prefix; padding is an
implementation device applied only after this prefix is formed.  For every
\(r\in\mathcal R_g\), define its nonempty input-cell set as
\begin{equation}
    \mathcal O_{g,r}
    =
    \left\{
        n\in\mathcal C_g^{\mathrm{in}}
        \;\middle|\;
        r(n)=r
    \right\}.
    \label{eq:row-input-cells}
\end{equation}
In particular,
\begin{equation}
    n_g^\star
    \notin
    \bigcup_{r\in\mathcal R_g}\mathcal O_{g,r}.
    \label{eq:target-exclusion}
\end{equation}
The database graph in Eq.~\eqref{eq:database-graph} is used only to select
the context.  Message passing is performed on the two query-specific graphs
defined next.


\subsection{Query-Specific Graph Construction}
\label{sec:graph-construction}

\subsubsection{Candidate-Augmented Row Graph}
\label{sec:main-graph}

For instance \(g\), we define the main prediction graph as a directed,
typed multigraph
\begin{equation}
    \mathcal G_g^{\main}
    =
    \bigl(
        \mathcal V_g^{\main},
        \mathcal E_g^{\main},
        \tau_g,
        \alpha_g;
        r_g^\star
    \bigr).
    \label{eq:main-graph}
\end{equation}
The semicolon emphasizes that \(r_g^\star\) is a distinguished query node,
not an additional node type.  The node set is the disjoint union
\begin{equation}
    \mathcal V_g^{\main}
    =
    \mathcal V_g^{\mathrm{row}}
    \mathbin{\dot\cup}
    \mathcal V_g^{\mathrm{cand}},
    \qquad
    \mathcal V_g^{\mathrm{row}}=\mathcal R_g,
    \label{eq:main-node-set}
\end{equation}
where the two binary candidate nodes are
\begin{equation}
    \mathcal V_g^{\mathrm{cand}}
    =
    \bigl\{
        v_{g,0}^{\mathrm{cand}},
        v_{g,1}^{\mathrm{cand}}
    \bigr\}.
    \label{eq:binary-candidates}
\end{equation}
The nodes \(v_{g,0}^{\mathrm{cand}}\) and
\(v_{g,1}^{\mathrm{cand}}\) represent the hypotheses
\textsc{False} and \textsc{True}, respectively.

The directed edge-index set is
\begin{equation}
    \mathcal E_g^{\main}
    =
    \mathcal E_g^{\mathsf{f2p}}
    \mathbin{\dot\cup}
    \mathcal E_g^{\mathsf{p2f}}
    \mathbin{\dot\cup}
    \mathcal E_g^{\mathsf{label}}.
    \label{eq:main-edge-set}
\end{equation}
The map
\begin{equation}
    \tau_g:
    \mathcal E_g^{\main}
    \longrightarrow
    \mathcal T_{\main},
    \qquad
    \mathcal T_{\main}
    =
    \{\mathsf{f2p},\mathsf{p2f},\mathsf{label}\},
    \label{eq:main-edge-type-map}
\end{equation}
assigns a type \(\tau_e\equiv\tau_g(e)\) to each edge.  The map
\(\alpha_g:\mathcal E_g^{\main}\to\R_{\geq0}\) assigns its scalar weight
\(\alpha_e\equiv\alpha_g(e)\).  All main-graph edges have
\(\alpha_e=1\) in the binary-classification model.  We retain the weight
notation because it makes the propagation operator compatible with future
soft-edge extensions.

\subsubsection{Foreign-Key and Observed-Label Edges}
\label{sec:main-edges}

The implementation stores a fixed-width FK-parent list and requires each
row to have at most five foreign-to-primary neighbors.  Subject to this
preprocessing constraint, the foreign-to-primary edge set is the restriction
of the stored database FK edges to the sampled rows:
\begin{equation}
    \mathcal E_g^{\mathsf{f2p}}
    =
    \left\{
        e\in\mathcal E^{\mathrm{FK}}
        \;\middle|\;
        \src(e),\dst(e)\in\mathcal R_g
    \right\}.
    \label{eq:f2p-edges}
\end{equation}
Thus, \(\src(e)\) is the sampled child row and \(\dst(e)\) is the sampled
parent row.  Every retained edge \(e\) is also represented by a distinct
reverse edge \(\bar e\):
\begin{equation}
    \mathcal E_g^{\mathsf{p2f}}
    =
    \left\{
        \bar e
        \;\middle|\;
        \begin{array}{l}
        e\in\mathcal E_g^{\mathsf{f2p}},\\[-1mm]
        \src(\bar e)=\dst(e),\qquad
        \dst(\bar e)=\src(e)
        \end{array}
    \right\}.
    \label{eq:p2f-edges}
\end{equation}
The two directions remain distinct edge types and therefore use distinct
message transformations.

Observed-label edges inject target-aligned contextual evidence into the
candidate nodes.  Define the set of context rows with a visible label as
\begin{equation}
    \mathcal R_g^{\mathrm{lab}}
    =
    \left\{
        r\in\mathcal R_g\setminus\{r_g^\star\}
        \;\middle|\;
        \begin{array}{l}
        \table(r)=\table(r_g^\star),\\[-1mm]
        \exists\,n\in\mathcal O_{g,r}:c(n)=c_g^\star
        \end{array}
    \right\}.
    \label{eq:observed-label-rows}
\end{equation}
For \(r\in\mathcal R_g^{\mathrm{lab}}\), let
\(y_{g,r}\in\{0,1\}\) be its observed
binary label.  We add the one-way edge
\begin{equation}
    \mathcal E_g^{\mathsf{label}}
    =
    \left\{
        \bigl(r,v_{g,y_{g,r}}^{\mathrm{cand}}\bigr)
        \;\middle|\;
        r\in\mathcal R_g^{\mathrm{lab}}
    \right\}.
    \label{eq:label-edges}
\end{equation}
No reverse candidate-to-row edge is added.  Since the seed is excluded from
\(\mathcal R_g^{\mathrm{lab}}\), its unknown target cannot create a label
edge.

\subsubsection{Auxiliary Relation Graph}
\label{sec:relation-graph-construction}

The main graph distinguishes only the generic directions
\(\mathsf{f2p}\) and \(\mathsf{p2f}\).  To represent the schema-level role
of an FK link, we construct an auxiliary directed, typed multigraph
\begin{equation}
    \mathcal G_g^{\rel}
    =
    \bigl(
        \mathcal V_g^{\rel},
        \mathcal E_g^{\rel},
        \eta_g;
        \mathcal Q_g^{\rel}
    \bigr).
    \label{eq:relation-graph}
\end{equation}
For each \(e\in\mathcal E_g^{\mathsf{f2p}}\), define its ordered
table-pair relation by
\begin{equation}
    \chi_g(e)
    =
    \bigl(\table(\src(e)),\table(\dst(e))\bigr),
    \label{eq:fk-to-relation}
\end{equation}
where \(\src(e)\) belongs to the child table and \(\dst(e)\) belongs to the
parent table.  The relation-node set contains the unique table pairs present in
the sampled graph:
\begin{equation}
    \mathcal V_g^{\rel}
    =
    \bigl\{
        \chi_g(e)
        \;\big|\;
        e\in\mathcal E_g^{\mathsf{f2p}}
    \bigr\}.
    \label{eq:relation-node-set}
\end{equation}
Multiple row-level FK links with the same ordered table pair therefore share
one relation node.  Relation nodes are local to query graph \(g\), and the
current construction does not distinguish multiple FK columns connecting
the same ordered pair of tables.  For a reverse edge in
\(\mathcal E_g^{\mathsf{p2f}}\), we formally set
\begin{equation}
    \chi_g(\bar e)
    \coloneqq
    \chi_g(e),
    \qquad
    e\in\mathcal E_g^{\mathsf{f2p}},
    \label{eq:reverse-fk-relation}
\end{equation}
so both directions map to the same relation node.

For a relation node \(\rho\in\mathcal V_g^{\rel}\), let
\(\operatorname{child}(\rho)\) and \(\operatorname{parent}(\rho)\) denote
the two components of its ordered table pair.  For distinct nodes
\(\rho,\rho'\in\mathcal V_g^{\rel}\), we define four directed structural
edge types:
\begin{align}
    (\rho,\rho',\mathsf{h2h})\in\mathcal E_g^{\mathsf{h2h}}
    &\iff
    \operatorname{child}(\rho)
    =
    \operatorname{child}(\rho'),
    \label{eq:h2h-edge}\\
    (\rho,\rho',\mathsf{t2t})\in\mathcal E_g^{\mathsf{t2t}}
    &\iff
    \operatorname{parent}(\rho)
    =
    \operatorname{parent}(\rho'),
    \label{eq:t2t-edge}\\
    (\rho,\rho',\mathsf{h2t})\in\mathcal E_g^{\mathsf{h2t}}
    &\iff
    \operatorname{child}(\rho)
    =
    \operatorname{parent}(\rho'),
    \label{eq:h2t-edge}\\
    (\rho,\rho',\mathsf{t2h})\in\mathcal E_g^{\mathsf{t2h}}
    &\iff
    \operatorname{parent}(\rho)
    =
    \operatorname{child}(\rho').
    \label{eq:t2h-edge}
\end{align}
Here \emph{head} corresponds to the child table and \emph{tail} to the
parent table.  The complete relation-edge set is
\begin{equation}
    \mathcal E_g^{\rel}
    =
    \mathcal E_g^{\mathsf{h2h}}
    \mathbin{\dot\cup}
    \mathcal E_g^{\mathsf{t2t}}
    \mathbin{\dot\cup}
    \mathcal E_g^{\mathsf{h2t}}
    \mathbin{\dot\cup}
    \mathcal E_g^{\mathsf{t2h}},
    \label{eq:relation-edge-set}
\end{equation}
and
\begin{equation}
    \eta_g:
    \mathcal E_g^{\rel}
    \longrightarrow
    \mathcal T_{\rel},
    \qquad
    \mathcal T_{\rel}
    =
    \{\mathsf{h2h},\mathsf{t2t},\mathsf{h2t},\mathsf{t2h}\}
    \label{eq:relation-edge-type-map}
\end{equation}
assigns each relation edge its type.  All relation-edge weights equal one.
The third tuple component is an explicit type tag.  Hence, if two
structural predicates hold for the same ordered endpoint pair, the
implementation retains parallel relation edges with different types.

\subsubsection{Query-Relation Marking}
\label{sec:query-relation-marking}

The relation graph is made specific to the current prediction by marking
the relations in which the seed participates as the child/source row.  The
query-relation set is
\begin{equation}
    \mathcal Q_g^{\rel}
    =
    \left\{
        \rho\in\mathcal V_g^{\rel}
        \;\middle|\;
        \exists\,e\in\mathcal E_g^{\mathsf{f2p}}
        \text{ such that }
        \src(e)=r_g^\star
        \text{ and }
        \chi_g(e)=\rho
    \right\}.
    \label{eq:query-relation-set}
\end{equation}
Thus, a relation is not marked merely because the seed is its parent.  This
set provides the query boundary used to initialize relation-node states in
Section~\ref{sec:relation-node-init}.


\subsection{Initial Node Representations}
\label{sec:node-initialization}

\subsubsection{Type-Aware Cell Encoding}
\label{sec:cell-encoding}

Let \(n\in\mathcal O_{g,r}\) index an observed cell in row \(r\), and let
\begin{equation}
    \sigma_{r,n}
    \in
    \Sigma_{\mathrm{sem}}
    =
    \{\mathsf{num},\mathsf{text},\mathsf{datetime},\mathsf{bool}\}
    \label{eq:semantic-types}
\end{equation}
be its semantic type.  The cell provides a value feature
\(\mathbf u_{r,n}^{\mathrm{val}}\in\R^{d_{\sigma_{r,n}}}\) and a
precomputed column-name feature
\(\mathbf u_{r,n}^{\mathrm{col}}\in\R^{d_{\mathrm{text}}}\), where
\(d_\sigma=d_{\mathrm{text}}\) for \(\sigma=\mathsf{text}\), and
\(d_\sigma=1\) for
\(\sigma\in\{\mathsf{num},\mathsf{datetime},\mathsf{bool}\}\).

We use a separate affine value encoder for each semantic type and a shared
column-name encoder.  The resulting cell representation is
\begin{equation}
    \boxed{
    \mathbf x_{r,n}
    =
    \RMSNorm_{\mathrm{col}}
    \left(
        \mathbf W_{\mathrm{col}}
        \mathbf u_{r,n}^{\mathrm{col}}
        +
        \mathbf b_{\mathrm{col}}
    \right)
    +
    \RMSNorm_{\sigma_{r,n}}
    \left(
        \mathbf W_{\mathrm{val},\sigma_{r,n}}
        \mathbf u_{r,n}^{\mathrm{val}}
        +
        \mathbf b_{\mathrm{val},\sigma_{r,n}}
    \right)
    }
    \in\R^d,
    \label{eq:cell-encoding}
\end{equation}
where
\begin{equation}
    \mathbf W_{\mathrm{col}}
    \in\R^{d\times d_{\mathrm{text}}},
    \quad
    \mathbf b_{\mathrm{col}}\in\R^d,
    \quad
    \mathbf W_{\mathrm{val},\sigma}
    \in\R^{d\times d_\sigma},
    \quad
    \mathbf b_{\mathrm{val},\sigma}\in\R^d.
    \label{eq:cell-parameter-domains}
\end{equation}
For an input \(\boldsymbol\zeta\in\R^d\), each RMS normalization module
indexed by \(\nu\in\{\mathrm{col}\}\cup\Sigma_{\mathrm{sem}}\) has the form
\begin{equation}
    \RMSNorm_{\nu}(\boldsymbol\zeta)
    =
    \boldsymbol\gamma_\nu
    \odot
    \frac{\boldsymbol\zeta}
    {\sqrt{d^{-1}\|\boldsymbol\zeta\|_2^2+\varepsilon_{\mathrm{rms}}}},
    \qquad
    \boldsymbol\gamma_\nu\in\R^d,
    \label{eq:rmsnorm}
\end{equation}
where \(\varepsilon_{\mathrm{rms}}>0\) is a numerical stabilizer and
\(\odot\) denotes element-wise multiplication.  Separate encoders and
normalizers accommodate heterogeneous value modalities, while the
column-name term distinguishes semantically different fields with similar
values.

\subsubsection{Permutation-Invariant Cell-to-Row Pooling}
\label{sec:row-pooling}

For a sampled row \(r\), define its mean-pooled representation by
\begin{equation}
    \overline{\mathbf x}_{g,r}
    =
    \frac{1}{|\mathcal O_{g,r}|}
    \sum_{n\in\mathcal O_{g,r}}
    \mathbf x_{r,n}
    \in\R^d,
    \label{eq:mean-pool}
\end{equation}
and its coordinate-wise maximum by
\begin{equation}
    [\mathbf x_{g,r}^{\max}]_q
    =
    \max_{n\in\mathcal O_{g,r}}
    [\mathbf x_{r,n}]_q,
    \qquad
    q=1,\ldots,d.
    \label{eq:max-pool}
\end{equation}
The preliminary row representation is
\begin{equation}
    \boxed{
    \overline{\mathbf h}_{g,r}
    =
    \mathbf W_{\mathrm{c2r}}
    \left[
        \overline{\mathbf x}_{g,r}
        \concat
        \mathbf x_{g,r}^{\max}
    \right]
    +
    \mathbf b_{\mathrm{c2r}}
    }
    \in\R^d,
    \label{eq:cell-to-row}
\end{equation}
where
\begin{equation}
    \mathbf W_{\mathrm{c2r}}
    \in\R^{d\times2d},
    \qquad
    \mathbf b_{\mathrm{c2r}}
    \in\R^d.
    \label{eq:c2r-domains}
\end{equation}
Mean pooling captures the overall row content, whereas max pooling preserves
salient cell-level signals.  Both operations are invariant to the order in
which a row's cells are serialized.

\subsubsection{Relative-Time and Seed Encoding}
\label{sec:time-seed-encoding}

Let \(T_r\) be the timestamp associated with row \(r\), measured in seconds.
We express time relative to the seed in days:
\begin{equation}
    \Delta t_{g,r}
    =
    \frac{T_{r_g^\star}-T_r}{86400}.
    \label{eq:relative-time}
\end{equation}
If either timestamp is unavailable, the implementation sets
\(\Delta t_{g,r}=0\).  Let the fixed wavelengths and angular frequencies be
\begin{equation}
    \lambda_f
    =
    10^{\frac{4.56f}{F-1}}
    \ \text{days},
    \qquad
    \omega_f
    =
    \frac{2\pi}{\lambda_f},
    \qquad
    f=0,\ldots,F-1.
    \label{eq:temporal-frequencies}
\end{equation}
Writing \(\boldsymbol\omega=[\omega_0,\ldots,\omega_{F-1}]^\top\in\R^F\),
the multi-scale Fourier feature is
\begin{equation}
    \boldsymbol\phi_{\mathrm{time}}(\Delta t_{g,r})
    =
    \left[
        \cos(\Delta t_{g,r}\boldsymbol\omega)
        \concat
        \sin(\Delta t_{g,r}\boldsymbol\omega)
    \right]
    \in\R^{2F},
    \label{eq:time-feature}
\end{equation}
where the trigonometric functions act element-wise and
\(\boldsymbol\omega\) is fixed rather than learned.  The initial state of a
row node is
\begin{equation}
    \boxed{
    \mathbf h_{g,r}^{(0)}
    =
    \overline{\mathbf h}_{g,r}
    +
    \mathbf W_{\mathrm{time}}
    \boldsymbol\phi_{\mathrm{time}}(\Delta t_{g,r})
    +
    \mathbf b_{\mathrm{time}}
    +
    \ind[r=r_g^\star]\mathbf e_{\mathrm{seed}}
    }
    \in\R^d,
    \label{eq:row-initialization}
\end{equation}
where
\begin{equation}
    \mathbf W_{\mathrm{time}}\in\R^{d\times2F},
    \qquad
    \mathbf b_{\mathrm{time}}\in\R^d,
    \qquad
    \mathbf e_{\mathrm{seed}}\in\R^d.
    \label{eq:time-seed-domains}
\end{equation}
The Fourier features provide a smooth representation from daily to
approximately century-scale horizons: short wavelengths resolve nearby
times, whereas long wavelengths vary slowly over long intervals.  The
shared seed marker identifies the current query boundary to the downstream
GNN.

\subsubsection{Binary Candidate Initialization}
\label{sec:candidate-init}

The False and True candidates are represented by the fixed scalar values
\(0\) and \(1\), respectively.  They are encoded using the same Boolean
value encoder and RMS normalization as Boolean database cells:
\begin{equation}
    \boxed{
    \mathbf h_{g,k}^{(0)}
    \equiv
    \mathbf h_{v_{g,k}^{\mathrm{cand}}}^{(0)}
    =
    \RMSNorm_{\mathsf{bool}}
    \left(
        \mathbf W_{\mathrm{val},\mathsf{bool}}k
        +
        \mathbf b_{\mathrm{val},\mathsf{bool}}
    \right)
    +
    \mathbf e_{\mathrm{cand}},
    \quad
    k\in\{0,1\}
    }
    \label{eq:candidate-initialization}
\end{equation}
with \(\mathbf e_{\mathrm{cand}}\in\R^d\).  Thus, False and True are
runtime nodes rather than two independently learned embedding parameters;
their distinction is induced by their fixed values under a shared encoder.
The vector \(\mathbf e_{\mathrm{cand}}\) is a learned node-type marker shared
by all candidate nodes and all tasks.

\subsubsection{Relation-Node Initialization}
\label{sec:relation-node-init}

For \(\rho\in\mathcal V_g^{\rel}\), the initial relation-node state is
\begin{equation}
    \boxed{
    \mathbf z_{g,\rho}^{(0)}
    =
    \ind[\rho\in\mathcal Q_g^{\rel}]
    \mathbf e_{\mathrm{relquery}}
    }
    \in\R^d,
    \label{eq:relation-initialization}
\end{equation}
where \(\mathbf e_{\mathrm{relquery}}\in\R^d\) is a learned query-relation
marker shared across graphs and tasks.  Non-query relation nodes are
initialized at zero; no schema-specific relation-ID embeddings are used.


\subsection{Relation-Conditioned Message Passing}
\label{sec:message-passing}

\subsubsection{Relation-Graph Propagation}
\label{sec:relation-propagation}

The auxiliary graph is processed before the main graph.  Let
\(\epsilon:\rho'\to\rho\) be a relation edge with type
\(\eta_\epsilon\equiv\eta_g(\epsilon)\in\mathcal T_{\rel}\).  At relation
layer \(p\in\{1,\ldots,P\}\), its message is
\begin{equation}
    \widetilde{\mathbf m}_{\epsilon}^{(p)}
    =
    \ReLU
    \left(
        \mathbf U_{\eta_\epsilon}^{(p)}
        \mathbf z_{g,\rho'}^{(p-1)}
        +
        \mathbf b_{\rel,\eta_\epsilon}^{(p)}
    \right)
    \in\R^d,
    \label{eq:relation-message}
\end{equation}
where
\begin{equation}
    \mathbf U_{\eta}^{(p)}\in\R^{d\times d},
    \qquad
    \mathbf b_{\rel,\eta}^{(p)}\in\R^d.
    \label{eq:relation-edge-parameter-domains}
\end{equation}
For the incoming relation-edge set
\begin{equation}
    \mathcal E_{g,\rel}^{-}(\rho)
    =
    \{\epsilon\in\mathcal E_g^{\rel}:\dst(\epsilon)=\rho\},
    \label{eq:relation-incoming-edges}
\end{equation}
we compute
\begin{equation}
    \widetilde{\mathbf a}_{g,\rho}^{(p)}
    =
    \frac{
        \displaystyle
        \sum_{\epsilon\in\mathcal E_{g,\rel}^{-}(\rho)}
        \widetilde{\mathbf m}_{\epsilon}^{(p)}
    }{
        \displaystyle
        \max\bigl(1,|\mathcal E_{g,\rel}^{-}(\rho)|\bigr)
    }
    \in\R^d.
    \label{eq:relation-aggregation}
\end{equation}
The relation-node update is
\begin{equation}
    \boxed{
    \mathbf z_{g,\rho}^{(p)}
    =
    \LayerNorm_{\rel}^{(p)}
    \left(
        \mathbf z_{g,\rho}^{(p-1)}
        +
        \ReLU
        \left(
            \mathbf U_{\mathrm{self}}^{(p)}
            \mathbf z_{g,\rho}^{(p-1)}
            +
            \mathbf b_{\rel,\mathrm{self}}^{(p)}
        \right)
        +
        \widetilde{\mathbf a}_{g,\rho}^{(p)}
    \right)
    }
    \label{eq:relation-update}
\end{equation}
with
\begin{equation}
    \mathbf U_{\mathrm{self}}^{(p)}\in\R^{d\times d},
    \qquad
    \mathbf b_{\rel,\mathrm{self}}^{(p)}\in\R^d.
    \label{eq:relation-self-domains}
\end{equation}
Each \(\LayerNorm_{\rel}^{(p)}\) has its own learned scale and shift,
\(\boldsymbol\gamma_{\rel}^{(p)},\boldsymbol\beta_{\rel}^{(p)}\in\R^d\).
Both propagation networks use the LayerNorm stabilizer
\(\varepsilon_{\mathrm{ln}}=10^{-5}\).  Relation layers and relation-edge
types do not share parameters.

\subsubsection{FK Relation Conditioning}
\label{sec:fk-conditioning}

After \(P=2\) relation layers, the condition attached to a directed main
graph FK edge \(e\in\mathcal E_g^{\mathsf{f2p}}\cup
\mathcal E_g^{\mathsf{p2f}}\) is
\begin{equation}
    \boxed{
    \boldsymbol\delta_e
    =
    \mathbf W_{\mathrm{cond}}
    \mathbf z_{g,\chi_g(e)}^{(P)}
    +
    \mathbf b_{\mathrm{cond}}
    }
    \in\R^d,
    \label{eq:fk-condition}
\end{equation}
where
\begin{equation}
    \mathbf W_{\mathrm{cond}}\in\R^{d\times d},
    \qquad
    \mathbf b_{\mathrm{cond}}\in\R^d.
    \label{eq:condition-domains}
\end{equation}
The forward and reverse directions of the same FK link satisfy
\(\boldsymbol\delta_{\bar e}=\boldsymbol\delta_e\), but they subsequently
use different source-node states and different edge-type transforms.  The
conditions are computed once and reused in every main-graph layer.  The
condition projection is initialized with zero weights and zero bias, so
relation conditioning is
initially inactive and is introduced gradually through learning.  We set
\(\boldsymbol\delta_e=\mathbf0\) for every observed-label edge.

\subsubsection{Relation-Conditioned Row-Graph Propagation}
\label{sec:main-propagation}

Let \(e:j\to i\) be a directed main-graph edge.  At layer
\(\ell\in\{1,\ldots,L\}\), its weighted typed message is
\begin{equation}
    \boxed{
    \mathbf m_{g,e}^{(\ell)}
    =
    \alpha_e
    \ReLU
    \left(
        \mathbf W_{\tau_e}^{(\ell)}
        \mathbf h_{g,j}^{(\ell-1)}
        +
        \mathbf b_{\tau_e}^{(\ell)}
        +
        \boldsymbol\delta_e
    \right)
    }
    \in\R^d,
    \label{eq:main-message}
\end{equation}
where
\begin{equation}
    \mathbf W_{\tau}^{(\ell)}\in\R^{d\times d},
    \qquad
    \mathbf b_{\tau}^{(\ell)}\in\R^d,
    \qquad
    \tau\in\mathcal T_{\main}.
    \label{eq:main-edge-domains}
\end{equation}
The incoming edge set and weighted degree of node \(i\) are
\begin{equation}
    \mathcal E_{g,\main}^{-}(i)
    =
    \{e\in\mathcal E_g^{\main}:\dst(e)=i\},
    \qquad
    \deg_{\alpha}(i)
    =
    \sum_{e\in\mathcal E_{g,\main}^{-}(i)}\alpha_e.
    \label{eq:main-incoming-degree}
\end{equation}
We use the clamped degree-normalized aggregate
\begin{equation}
    \boxed{
    \mathbf a_{g,i}^{(\ell)}
    =
    \frac{
        \displaystyle
        \sum_{e\in\mathcal E_{g,\main}^{-}(i)}
        \mathbf m_{g,e}^{(\ell)}
    }{
        \displaystyle
        \max\bigl(1,\deg_{\alpha}(i)\bigr)
    }
    }
    \in\R^d.
    \label{eq:main-aggregation}
\end{equation}
The lower bound in the denominator prevents division by zero; consequently,
Eq.~\eqref{eq:main-aggregation} is not necessarily a strict weighted average
when the total incoming weight is below one.

The node update combines a residual path, a learned self transformation,
and the incoming aggregate:
\begin{equation}
    \boxed{
    \mathbf h_{g,i}^{(\ell)}
    =
    \LayerNorm_{\main}^{(\ell)}
    \left(
        \mathbf h_{g,i}^{(\ell-1)}
        +
        \ReLU
        \left(
            \mathbf W_{\mathrm{self}}^{(\ell)}
            \mathbf h_{g,i}^{(\ell-1)}
            +
            \mathbf b_{\mathrm{self}}^{(\ell)}
        \right)
        +
        \mathbf a_{g,i}^{(\ell)}
    \right)
    }
    \in\R^d,
    \label{eq:main-update}
\end{equation}
where
\begin{equation}
    \mathbf W_{\mathrm{self}}^{(\ell)}\in\R^{d\times d},
    \qquad
    \mathbf b_{\mathrm{self}}^{(\ell)}\in\R^d.
    \label{eq:main-self-domains}
\end{equation}
Each \(\LayerNorm_{\main}^{(\ell)}\) has its own learned scale and shift,
\(\boldsymbol\gamma_{\main}^{(\ell)},
\boldsymbol\beta_{\main}^{(\ell)}\in\R^d\).  The initial main-graph state
stacks all row and candidate states:
\begin{equation}
    \mathbf H_g^{(0)}
    =
    \operatorname{Concat}_{\mathrm{node}}
    \left(
        \{\mathbf h_{g,r}^{(0)}:r\in\mathcal R_g\},
        \{\mathbf h_{g,k}^{(0)}:k\in\{0,1\}\}
    \right)
    \in\R^{(|\mathcal R_g|+2)\times d}.
    \label{eq:main-initial-matrix}
\end{equation}
The concatenation uses a fixed ordering with row nodes first and candidate
nodes second.
The model then applies \(L=6\) independently parameterized layers:
\begin{equation}
    \mathbf H_g^{(\ell)}
    =
    \operatorname{MainMPLayer}^{(\ell)}
    \left(
        \mathbf H_g^{(\ell-1)},
        \mathcal G_g^{\main},
        \{\boldsymbol\delta_e\}
    \right),
    \qquad
    \ell=1,\ldots,L.
    \label{eq:main-layer-stack}
\end{equation}
Within a layer, the three edge types use distinct parameter sets.  Across
layers, no message or self-transformation parameters are shared.  All these
parameters are nevertheless shared across query graphs and prediction
tasks.  Type-specific transforms preserve the asymmetric semantics of
child-to-parent, parent-to-child, and observed-label evidence.  Degree
normalization controls the scale of accumulated evidence, while the
residual and self-transformation paths preserve and refine a node's own
state even when it has no incoming edges.

\subsubsection{Base Model as a Special Case}
\label{sec:base-special-case}

The base model and the relation-conditioned model share the same main graph,
node initialization, message-passing architecture, and readout.  They differ
only in the FK condition:
\begin{equation}
    \boxed{
    \boldsymbol\delta_e
    =
    \begin{cases}
        \mathbf0,
        &
        \text{base model or label edge},\\[1mm]
        \mathbf W_{\mathrm{cond}}
        \mathbf z_{g,\chi_g(e)}^{(P)}
        +
        \mathbf b_{\mathrm{cond}},
        &
        \text{relation-conditioned FK edge}.
    \end{cases}
    }
    \label{eq:base-rel-condition}
\end{equation}
Thus, the base model is recovered exactly by setting all FK conditions to
zero.  Importantly, the relation condition is added \emph{after} the
edge-type affine transformation, as in
\(\mathbf W_{\tau_e}^{(\ell)}\mathbf h_{g,j}^{(\ell-1)}
+\mathbf b_{\tau_e}^{(\ell)}+\boldsymbol\delta_e\); it is not transformed
together with the source state.


\subsection{Candidate Scoring and Classification Objective}
\label{sec:classification-readout}

After \(L\) main-graph layers, let
\(\mathbf h_{g,r_g^\star}^{(L)}\in\R^d\) be the final seed representation and
let \(\mathbf h_{g,k}^{(L)}\in\R^d\) be the final representation of
candidate \(k\in\{0,1\}\).  We concatenate them to form
\begin{equation}
    \boldsymbol\xi_{g,k}^{\mathrm{ro}}
    =
    \mathbf h_{g,r_g^\star}^{(L)}
    \concat
    \mathbf h_{g,k}^{(L)}
    \in\R^{2d}.
    \label{eq:readout-input}
\end{equation}
A two-layer MLP shared by both candidates and all tasks computes
\begin{align}
    \boldsymbol\psi_{g,k}^{\mathrm{ro}}
    &=
    \ReLU
    \left(
        \mathbf W_{\mathrm{ro},1}
        \boldsymbol\xi_{g,k}^{\mathrm{ro}}
        +
        \mathbf b_{\mathrm{ro},1}
    \right)
    \in\R^d,
    \label{eq:readout-hidden}\\
    s_{g,k}
    &=
    \mathbf w_{\mathrm{ro},2}^{\top}
    \boldsymbol\psi_{g,k}^{\mathrm{ro}}
    +
    b_{\mathrm{ro},2}
    \in\R,
    \label{eq:candidate-score}
\end{align}
where
\begin{equation}
    \mathbf W_{\mathrm{ro},1}\in\R^{d\times2d},
    \qquad
    \mathbf b_{\mathrm{ro},1}\in\R^d,
    \qquad
    \mathbf w_{\mathrm{ro},2}\in\R^d,
    \qquad
    b_{\mathrm{ro},2}\in\R.
    \label{eq:readout-domains}
\end{equation}
The scalar \(s_{g,0}\) scores the False hypothesis and \(s_{g,1}\) scores
the True hypothesis.  Let \(\mathcal X_g\) denote the complete constructed
model input: the main graph together with \(\{\boldsymbol\delta_e\}\), whose
values are either zero in the base model or produced by the relation graph.
We define the binary logit and probability by
\begin{equation}
    \Delta_g
    =
    s_{g,1}-s_{g,0},
    \qquad
    \pi_g
    =
    \sigmoid(\Delta_g)
    =
    \Pr\!\left(y_g=1\mid\mathcal X_g\right).
    \label{eq:classification-probability}
\end{equation}
The corresponding hard prediction is
\begin{equation}
    \widehat y_g
    =
    \underset{k\in\{0,1\}}{\arg\max}\;s_{g,k},
    \label{eq:classification-prediction}
\end{equation}
and the per-instance binary cross-entropy loss is
\begin{equation}
    \mathcal L_g^{\mathrm{clf}}
    =
    -y_g\log \pi_g
    -(1-y_g)\log(1-\pi_g).
    \label{eq:classification-loss}
\end{equation}
This classification-only formulation considers minibatches containing
binary queries.  For a minibatch \(\mathcal B\) containing
\(G=|\mathcal B|\) valid query instances, the training objective is
\begin{equation}
    \boxed{
    \mathcal L_{\mathcal B}
    =
    \frac{1}{G}
    \sum_{g\in\mathcal B}
    \mathcal L_g^{\mathrm{clf}}
    }.
    \label{eq:batch-loss}
\end{equation}


\subsection{Extension to Regression}
\label{sec:regression-placeholder}

\emph{Placeholder.  This subsection will specify continuous candidate
construction, the mapping from observed continuous labels to candidate
nodes, the regression training objective, continuous-valued inference, and
the combined classification--regression minibatch objective used by the
full implementation.}

\end{document}
